% ============================================================
% 新领域学习模板
% ============================================================

\section{新领域学习模板}

\subsection{领域入口页}

每进入一个新领域，先建立一页入口笔记，格式如下：

\begin{conceptbox}
\textbf{领域名称：} \underline{\hspace{8cm}}\\
\textbf{我为什么要学它：} \underline{\hspace{7cm}}\\
\textbf{最终要达到的能力：} \underline{\hspace{6.6cm}}\\
\textbf{时间盒：} \underline{\hspace{8.5cm}}
\end{conceptbox}

\subsection{问题地图}

\begin{center}
{\small
\begin{tabular}{p{3.2cm}p{4.8cm}p{5.2cm}}
\hline
\textbf{核心问题} & \textbf{为什么重要} & \textbf{代表方法 / 论文} \\
\hline
问题 1 & & \\
问题 2 & & \\
问题 3 & & \\
问题 4 & & \\
\hline
\end{tabular}
}
\end{center}

\subsection{方法对比表}

\begin{center}
{\small
\begin{tabular}{p{2.8cm}p{3.4cm}p{3.4cm}p{3.4cm}}
\hline
\textbf{方法} & \textbf{核心思想} & \textbf{解决的问题} & \textbf{代价 / 局限} \\
\hline
方法 A & & & \\
方法 B & & & \\
方法 C & & & \\
\hline
\end{tabular}
}
\end{center}

\subsection{每篇论文固定提取}

\begin{enumerate}[leftmargin=1.5em]
  \item \textbf{一句话贡献：} 这篇论文到底做了什么？
  \item \textbf{旧问题：} 它接在哪条技术路线后面？
  \item \textbf{关键机制：} 新的 loss、数据、模型结构或训练流程是什么？
  \item \textbf{实验结论：} 哪些结果真正支撑 claim？
  \item \textbf{我的判断：} 值不值得继续追？为什么？
\end{enumerate}

\subsection{掌握度检查}

\begin{summarybox}
一个领域可以暂时算“入门掌握”，当我能够完成下面五件事：
\begin{itemize}[leftmargin=1.5em]
  \item 画出主流方法谱系；
  \item 不看笔记解释 3 个核心算法；
  \item 写出至少 1 个关键损失函数，并解释每一项；
  \item 对一篇新论文做路线归类和贡献判断；
  \item 说出这个领域未来 3 个可能的研究方向。
\end{itemize}
\end{summarybox}
